\chapter{Introducción específica}
\label{Chapter2}

En este capítulo se describen las principales plataformas de hardware, componentes externos, capas de software y protocolos de comunicación utilizados durante el desarrollo del trabajo. Asimismo, se presentan las tecnologías de terceros empleadas para la implementación de las funciones de adquisición, procesamiento y comunicación del sistema.

%----------------------------------------------------------------------------------------

\section{Plataforma de hardware}

Tanto la estación central como las estaciones remotas del sistema utilizan el microcontrolador STM32H563RI de STMicroelectronics, perteneciente a la familia STM32H5 \cite{STM32H563,STMicroelectronics,STM32H5}. Esta decisión permitió uniformar la arquitectura de hardware y simplificar el desarrollo y mantenimiento del firmware. Además, facilitó la reutilización de controladores, configuraciones de periféricos y capas de abstracción de hardware entre ambos tipos de nodos.

El STM32H563RI se basa en una arquitectura ARM Cortex-M33 de 32 bits \cite{CortexM33}, orientada a sistemas embebidos de alto rendimiento y plataformas conectadas. El núcleo incorpora una unidad de punto flotante FPU (\textit{Floating Point Unit}, Unidad de Punto Flotante), extensiones DSP (\textit{Digital Signal Processing}, Procesamiento Digital de Señales) y mecanismos de seguridad basados en TrustZone \cite{STM32H5,STM32H563_RefManual}, adecuados para aplicaciones industriales e IoT.

El microcontrolador dispone de memoria Flash integrada y memoria SRAM (\textit{Static Random Access Memory}, Memoria Estática de Acceso Aleatorio) interna, además de una amplia variedad de periféricos de comunicación y control \cite{STM32H563}. Entre los periféricos utilizados durante el desarrollo se encuentran interfaces SPI (\textit{Serial Peripheral Interface}, Interfaz Periférica Serie), UART (\textit{Universal Asynchronous Receiver-Transmitter}, Transmisor-Receptor Asíncrono Universal) y Ethernet MAC (\textit{Media Access Control}, Control de Acceso al Medio).

La selección de este microcontrolador también se fundamentó en la disponibilidad de herramientas de desarrollo y middleware provistos por STMicroelectronics, incluyendo STM32CubeIDE \cite{stm32cubeide}, bibliotecas HAL (\textit{Hardware Abstraction Layer}, Capa de Abstracción de Hardware) \cite{hal} y soporte para sistemas operativos de tiempo real como FreeRTOS (\textit{Free Real-Time Operating System}, Sistema Operativo de Tiempo Real Libre) \cite{freertos}. Estas herramientas simplificaron significativamente el desarrollo, depuración y mantenimiento del firmware del sistema.

%----------------------------------------------------------------------------------------

\section{Componentes y módulos externos}

Complementariamente a la plataforma de procesamiento basada en el STM32H563RI, el sistema incorpora distintos circuitos integrados y módulos externos orientados a la adquisición de señales industriales, manejo de salidas, comunicación cableada e inalámbrica y validación del firmware durante las etapas de desarrollo. A continuación, se describen los componentes más relevantes utilizados en el trabajo.

\subsection{SCLT3-8BT8}

El SCLT3-8BT8 de STMicroelectronics es una interfaz de entradas digitales industriales de ocho canales con transferencia serializada de estados mediante interfaz SPI \cite{sclt38bt8}. El dispositivo está orientado a aplicaciones de automatización industrial y PLC. Además, proporciona adaptación y protección para señales de 24 V provenientes de sensores y dispositivos industriales.

El circuito incorpora protección contra sobretensiones y compatibilidad con estándares industriales de inmunidad electromagnética, incluyendo IEC61000-4-2, IEC61000-4-4 e IEC61000-4-5. Además, dispone de indicadores de estado integrados y mecanismos de diagnóstico para detección de fallas en las entradas.

\subsection{VNI8200XP-32}

El VNI8200XP-32 de STMicroelectronics es un controlador de salidas digitales industriales de ocho canales basado en transistores MOSFET (\textit{Metal-Oxide-Semiconductor Field-Effect Transistor}, Transistor de Efecto de Campo Metal-Óxido-Semiconductor) de tipo \textit{high-side} \cite{vni8200xp32}. El dispositivo permite accionar cargas industriales de 24 V mediante una interfaz SPI o paralelo configurable.

El componente incorpora funciones de protección y diagnóstico, incluyendo protección térmica, limitación de corriente, detección de sobrecorriente, monitoreo de alimentación y señalización de fallas. También dispone de mecanismos específicos frente a condiciones de cortocircuito y sobretemperatura, características relevantes en aplicaciones industriales.

\subsection{SN65HVD82}

El SN65HVD82 de Texas Instruments es un transceptor RS-485 (\textit{Recommended Standard 485}, Estándar de Comunicación Diferencial Serie) diseñado para comunicaciones industriales robustas \cite{sn65hvd82}. El dispositivo opera con alimentación de 5 V y proporciona transmisión diferencial compatible con buses multipunto de larga distancia.

El componente incorpora protección ESD (\textit{Electrostatic Discharge}, Descarga Electrostática) conforme a normas IEC orientadas a aplicaciones industriales y permite operación en entornos con elevado nivel de ruido electromagnético. Además, soporta altas velocidades de transmisión y presenta bajo consumo energético.

\subsection{MRF24J40}

El MRF24J40 de Microchip es un transceptor inalámbrico compatible con el estándar IEEE 802.15.4 \cite{mrf24j40}. El dispositivo opera en la banda ISM (\textit{Industrial, Scientific and Medical}, Industrial, Científica y Médica) de 2.4 GHz y está orientado a aplicaciones de redes inalámbricas de bajo consumo.

El componente incorpora las capas PHY y MAC, además de mecanismos de cifrado AES-128 (\textit{Advanced Encryption Standard}, Estándar de Cifrado Avanzado). La comunicación con el microcontrolador se realiza mediante interfaz SPI.

\subsection{LAN8742A}

El LAN8742A de Microchip es un transceptor Ethernet PHY compatible con Ethernet 10BASE-T y 100BASE-TX \cite{lan8742a}. El dispositivo se comunica con el microcontrolador mediante interfaz RMII (\textit{Reduced Media Independent Interface}, Interfaz Independiente del Medio Reducida), lo que permite implementar conectividad Ethernet de bajo número de pines.

El componente incorpora soporte para auto-negociación, detección automática de polaridad y tecnología HP Auto-MDIX (\textit{Automatic Medium-Dependent Interface Crossover}, Cruce Automático de Interfaz Dependiente del Medio), característica que permite utilizar cables directos o cruzados. Además, dispone de mecanismos de administración de energía y soporte para Wake-on-LAN (\textit{Wake on Local Area Network}, Activación mediante Red Local).

\subsection{Kits de desarrollo y evaluación}

Durante las etapas iniciales del proyecto se utilizaron placas STM32 NUCLEO-H563ZI \cite{nucleoh563}, basadas en el microcontrolador STM32H563ZI. Estas placas permitieron realizar pruebas funcionales y tareas de depuración temprana antes del desarrollo del hardware definitivo.

Asimismo, se empleó la placa de expansión X-NUCLEO-PLC01A1 \cite{nucleoplc}, orientada a aplicaciones PLC industriales. Esta placa integra los dispositivos SCLT3-8BT8 \cite{sclt38bt8} y VNI8200XP-32 \cite{vni8200xp32}, que proporcionan entradas y salidas industriales de 24 V manejadas mediante SPI. Su utilización permitió validar la arquitectura de firmware y los controladores desarrollados para las interfaces industriales utilizadas posteriormente en el diseño final.

%----------------------------------------------------------------------------------------

\section{Sistema operativo y middleware}

El desarrollo del firmware del sistema se realizó utilizando distintas capas de software orientadas a facilitar la abstracción del hardware, la administración de tareas concurrentes y la implementación de protocolos de comunicación.

\subsection{FreeRTOS}

FreeRTOS es un sistema operativo de tiempo real orientado a sistemas embebidos \cite{freertos}. El sistema proporciona mecanismos de planificación multitarea basados en prioridades, sincronización mediante semáforos y mutex, comunicación entre tareas y administración de temporizadores de software.

La utilización de un RTOS permitió implementar múltiples tareas concurrentes asociadas a adquisición de datos, comunicación inalámbrica y procesamiento de protocolos Ethernet. FreeRTOS también facilitó la modularización del firmware y la separación funcional entre los distintos subsistemas de la aplicación.

\subsection{HAL y BSP}

Para el acceso y configuración de los periféricos del microcontrolador se utilizaron las bibliotecas HAL (\textit{Hardware Abstraction Layer}, Capa de Abstracción de Hardware) y BSP (\textit{Board Support Package}, Paquete de Soporte de Placa) provistas por STMicroelectronics \cite{hal,bsp}.

La capa HAL proporciona funciones de alto nivel para el manejo de periféricos internos del microcontrolador, incluyendo interfaces SPI, UART, Ethernet y temporizadores. Esta capa simplifica el desarrollo del firmware mediante una interfaz uniforme e independiente de registros específicos del hardware.

Por otro lado, la capa BSP incorpora controladores y configuraciones específicas para placas de desarrollo y módulos de evaluación utilizados durante el proyecto, facilitando la inicialización y utilización de periféricos externos integrados en dichas plataformas.

\subsection{lwIP}

Para la implementación de conectividad Ethernet y protocolos TCP/IP (\textit{Transmission Control Protocol/Internet Protocol}, Protocolo de Control de Transmisión/Protocolo de Internet) se utilizó el \textit{stack} lwIP (\textit{Lightweight IP}) \cite{tcpip,lwip}.

lwIP es una implementación liviana del protocolo TCP/IP \cite{tcpip} orientada a sistemas embebidos con recursos limitados. El \textit{stack} incorpora soporte para protocolos IPv4, TCP, UDP (\textit{User Datagram Protocol}, Protocolo de Datagrama de Usuario), ICMP (\textit{Internet Control Message Protocol}, Protocolo de Mensajes de Control de Internet), DHCP (\textit{Dynamic Host Configuration Protocol}, Protocolo de Configuración Dinámica de Host) y ARP (\textit{Address Resolution Protocol}, Protocolo de Resolución de Direcciones).

En el sistema desarrollado, lwIP fue integrado junto con FreeRTOS y el periférico Ethernet del STM32H563RI. Esta integración permitió implementar comunicación Ethernet para supervisión, configuración y transmisión de datos mediante redes TCP/IP.

\section{Protocolos e interfaces de comunicación}

El sistema desarrollado utiliza distintas interfaces y protocolos de comunicación orientados a la adquisición de datos, control de periféricos, conectividad Ethernet y depuración del firmware. La selección de cada interfaz se realizó en función de los requerimientos de velocidad, complejidad de implementación, robustez y compatibilidad con los dispositivos externos utilizados en el proyecto.

\subsection{SPI}

La interfaz SPI fue utilizada para la comunicación entre el microcontrolador y distintos circuitos integrados externos, incluyendo las interfaces industriales SCLT3-8BT8 \cite{sclt38bt8} y VNI8200XP-32 \cite{vni8200xp32}, así como también el transceptor inalámbrico MRF24J40 \cite{mrf24j40}.

SPI es un protocolo de comunicación serial síncrono basado en una arquitectura maestro-esclavo, caracterizado por su elevada velocidad de transferencia y baja complejidad de implementación. La interfaz utiliza líneas dedicadas de transmisión, recepción y señal de reloj, además de señales independientes de selección de dispositivo.

\subsection{UART y RS-485}

La interfaz UART fue utilizada junto con el transceptor SN65HVD82 \cite{sn65hvd82} para implementar comunicación RS-485 (\textit{Recommended Standard 485}, Estándar de Comunicación Diferencial Serie) \cite{rs485}.

UART permite establecer comunicaciones seriales asíncronas mediante transmisión y recepción de datos sin necesidad de señal de reloj compartida. En conjunto con RS-485, posibilita implementar enlaces diferenciales robustos, adecuados para entornos industriales con elevado nivel de ruido electromagnético y largas distancias de cableado.

\subsection{Ethernet y RMII}

La conectividad Ethernet del sistema se implementó mediante el periférico Ethernet MAC integrado en el microcontrolador STM32H563RI \cite{STM32H563_RefManual} y el transceptor Ethernet PHY LAN8742A \cite{lan8742a}.

La comunicación entre el microcontrolador y el PHY se realizó mediante la interfaz RMII (\textit{Reduced Media Independent Interface}, Interfaz Independiente del Medio Reducida) \cite{rmii}, utilizada para transmisión y recepción de datos Ethernet con bajo número de pines.

La capa física Ethernet permite establecer comunicación 10BASE-T y 100BASE-TX sobre medio cableado de par trenzado. Esto facilita la integración con redes Ethernet industriales y servicios basados en TCP/IP.

\subsection{GPIO}

Las interfaces GPIO (\textit{General Purpose Input/Output}, Entrada/Salida de Propósito General) fueron utilizadas para el manejo de señales digitales de control, monitoreo de estados y accionamiento de indicadores luminosos.

Estas señales incluyen líneas de habilitación de periféricos, señales de interrupción, selección de dispositivos SPI, monitoreo de estados de alimentación y control de LEDs indicadores presentes en la placa.

\subsection{SWD}

La programación y depuración del microcontrolador se realizó mediante la interfaz SWD (\textit{Serial Wire Debug}, Depuración Serie de Dos Hilos) \cite{swi}, utilizada por las herramientas de desarrollo STMicroelectronics.

SWD constituye una interfaz de depuración derivada de JTAG (\textit{Joint Test Action Group}) \cite{jtag}. Esta interfaz reduce la cantidad de señales requeridas y conserva capacidades de programación, ejecución paso a paso, lectura de memoria y monitoreo interno del sistema embebido.